\documentclass[UTF8,fontset=fandol,11pt]{ctexart}

\usepackage[a4paper,margin=2.35cm]{geometry}
\usepackage{booktabs}
\usepackage{array}
\usepackage{tabularx}
\usepackage{graphicx}
\usepackage{float}
\usepackage{caption}
\usepackage{xcolor}
\usepackage{enumitem}
\usepackage{amsmath}
\usepackage{hyperref}

\definecolor{chapterblue}{HTML}{2F5F85}
\definecolor{softgray}{HTML}{F5F6F7}
\definecolor{textgray}{HTML}{333333}

\graphicspath{{../../charts/}}
\hypersetup{
  colorlinks=true,
  linkcolor=chapterblue,
  urlcolor=chapterblue,
  citecolor=chapterblue
}
\captionsetup{font=small,labelfont=bf}
\setlist[itemize]{leftmargin=1.6em,itemsep=0.15em}
\setlength{\parindent}{2em}
\setlength{\parskip}{0.35em}
\renewcommand{\arraystretch}{1.18}

\title{\textbf{第五部分：数据准备、描述性分析与建模交接}}
\author{DASE7337 Group Project: Airline Delay and Cancellation Operational Risk}
\date{April 2026}

\begin{document}
\maketitle

\begin{abstract}
本章节承接前文对 JFK 航班延误运营风险的行业背景、风险因素、SWOT 与 PEST 分析，进一步说明本项目如何把 BTS On-Time Performance 原始数据整理成可以直接服务于后续模型构建的数据资产。章节重点不在于单纯展示清洗结果，而是解释每一步处理背后的风险管理逻辑：为什么要统一字段口径，为什么要删除无实际运营量的记录，为什么要同时保留频率变量、严重度变量和原因分类变量，以及这些输出如何支撑后续的频率模型、严重度模型、状态划分与 aggregate risk 模拟。
\end{abstract}

\setcounter{section}{4}
\section{数据准备与描述性分析}

\subsection{章节定位}
前一章节已经说明，JFK 的航班延误风险既包含外部冲击，例如天气与空域约束，也包含内部运营链条中的排班、地面服务和飞机周转问题。因此，本章节的数据工作有一个明确目标：把原始运营记录转化为一组清晰、稳定、可复现的风险度量。换句话说，数据清洗不是技术性的附属步骤，而是后续建模能否成立的基础。

本项目使用的数据来自 BTS On-Time Performance 中的 Airline Delay Cause 数据表，样本限定为 2020 年 1 月至 2025 年 12 月 JFK 到达航班的航司月度记录。原始数据共 570 条记录，清洗后保留 569 条。最终数据覆盖 72 个机场月份、12 家航司、657,671 架次到达航班，并记录了 137,247 次 15 分钟以上到达延误、18,230 次取消、1,873 次改降或备降，以及 11,125,838 分钟到达延误。

\subsection{数据处理流程}
数据处理流程可以概括为四步。第一步是读取原始 CSV 与官方字段说明文件，确认字段与 BTS 定义一致。第二步是统一航司名称和机场名称，避免同一实体因为命名差异被拆成多个类别。第三步是把原始字段改写成更直观的变量名，例如将 \texttt{arr\_flights} 改为 \texttt{total\_arrival\_flights}，将 \texttt{arr\_del15} 改为 \texttt{delayed\_arrivals\_15\_plus}。第四步是建立核心清洗数据、字段字典、月度汇总、延误原因汇总和航司风险画像。

\begin{table}[H]
\centering
\caption{数据清洗与样本覆盖情况}
\begin{tabular}{lr}
\toprule
指标 & 数值 \\
\midrule
原始航司月份记录 & 570 \\
清洗后航司月份记录 & 569 \\
因总到达航班量缺失或为零而删除 & 1 \\
因延误航班数超过总航班数而删除 & 0 \\
覆盖年份 & 2020--2025 \\
机场月份数量 & 72 \\
航司数量 & 12 \\
总到达航班量 & 657,671 \\
总延误分钟 & 11,125,838 \\
\bottomrule
\end{tabular}
\end{table}

这里最重要的清洗规则有两条。第一，删除总到达航班量缺失或小于等于零的记录，因为这类记录无法形成有效的延误率、取消率或严重度指标。第二，检查延误航班数是否超过总到达航班量，这是一类明显的逻辑异常。当前样本没有出现第二类错误，但保留这一检查可以保证未来更新数据时流程仍然稳健。

\subsection{变量设计与模型口径}
清洗后的核心数据保留了三类变量。第一类是频率变量，例如到达航班量、15 分钟以上延误航班数、取消数和改降数。第二类是严重度变量，例如总到达延误分钟，以及 Airline、Weather、NAS、Security、Late Aircraft 五类原因对应的延误分钟。第三类是识别变量，包括年份、月份和航司名称。这一结构与后续 frequency--severity 建模框架直接对应。

在航司月份层面，最基本的频率口径可以写为
\[
\mathrm{DelayRate}_{a,t} =
\frac{\mathrm{DelayedArrivals}_{a,t}}{\mathrm{TotalArrivals}_{a,t}},
\qquad
\mathrm{CancellationRate}_{a,t} =
\frac{\mathrm{CancelledArrivals}_{a,t}}{\mathrm{TotalArrivals}_{a,t}}.
\]
严重度口径则用平均每次延误分钟数表示：
\[
\mathrm{AverageSeverity}_{a,t} =
\frac{\mathrm{TotalDelayMinutes}_{a,t}}{\mathrm{DelayedArrivals}_{a,t}}.
\]
当分母为零时，后续建模应把该航司月份视为没有可观测延误事件，而不是人为制造一个极端严重度数值。这样处理有助于把”是否发生事件”和”发生后损失多大”分开，正好对应运营风险中的频率与严重度分解。

\subsubsection*{Delay-Equivalent Minutes：统一损失口径}
由于 BTS 数据中取消和备降并不以延误分钟形式记录，项目将其转换为统一的 \emph{delay-equivalent minutes}（延误等价分钟）。某航司在月份 $t$ 的总运营冲击定义为
\[
L_{a,t} = \mathrm{TotalDelayMinutes}_{a,t}
         + 180 \times \mathrm{CancelledArrivals}_{a,t}
         + 240 \times \mathrm{DivertedArrivals}_{a,t}.
\]
取消航班赋予 180 分钟、备降航班赋予 240 分钟的选择，同时基于监管阈值、国际旅客权益惯例和备降恢复流程的实证研究。

\textbf{为什么取消航班的 penalty 是 180 分钟？}
美国交通部停机坪延误规则（14~CFR~259.4）自 2010 年 4 月 29 日起规定，国内航班在停机坪上的等待时间不得超过 3 小时（180 分钟），违规航司将面临每位旅客最高 27,500 美元的罚款 \cite{dot2010tarmac}。这一规定在监管层面划定了”严重延误”与”实质取消”之间的法定分界线：当延误接近 3 小时时，航司往往选择主动取消航班以规避违规风险。美国交通部监察长办公室（OIG）的审计报告证实，该规则实施后的前三年内航班取消率显著上升，说明 3 小时阈值确实是航司在延误与取消之间做出决策的实际临界点 \cite{dotoig2014tdr}。与此同时，欧盟旅客权益条例 EC~261/2004 将到达延误超过 3 小时的航班在赔偿义务上等同于取消航班处理，旅客可获得 250 至 600 欧元的固定赔偿 \cite{eu261}，进一步确认了 3 小时作为延误与取消等价边界的国际共识。从旅客实际影响看，航班取消后通常需要等待下一个可用航班，对于 JFK 这样的大型枢纽，主要航线间隔通常在 2--4 小时之间，因此 180 分钟是保守的下界估计。

\textbf{为什么备降航班的 penalty 是 240 分钟？}
同一法规对国际航班设定的停机坪等待上限为 4 小时（240 分钟）\cite{dot2010tarmac}。JFK 作为美国最重要的国际航空枢纽之一，其备降事件中有相当比例涉及国际航班，因此 240 分钟与该监管阈值直接对应。更重要的是，备降的运营冲击严格大于取消：飞机已经起飞、消耗燃油并降落在非计划机场，后续恢复流程包括备降机场的地面处理、可能的加油、旅客重新安检，以及继续飞往原目的地或安排地面交通。Peterson 等人（2021）对美国航班备降数据的研究表明，单次备降的平均直接运营成本约为 38,000 美元，且非医疗原因的备降恢复时间显著长于医疗原因备降 \cite{peterson2021diversion}。因此 4 小时是对备降全流程额外延误的合理下界估计。$240 > 180$ 的设定也保持了直觉上的严重度递进关系：备降比取消多出了”已在空中”的阶段，涉及额外的燃油消耗、机组工时、备降机场资源占用以及后续航班的连锁影响。

综上，180 分钟与 240 分钟的选择并非任意设定，而是同时锚定于美国联邦监管阈值（DOT Tarmac Delay Rule）、国际旅客权益惯例（EU~261/2004）以及备降恢复流程的实证研究。这一设计使得 delay-equivalent minutes 作为统一损失口径既有监管依据，又有运营逻辑支撑，同时避免了引入主观性更强的货币化假设。

\begin{thebibliography}{9}
\bibitem{dot2010tarmac}
U.S.\ Department of Transportation, ``Enhancing Airline Passenger Protections,'' 75~FR~32318, June~2010. 14~CFR~259.4.

\bibitem{dotoig2014tdr}
U.S.\ DOT Office of Inspector General, ``Effects of the Tarmac Delay Rule on Flight Cancellations and Delays,'' Report No.\ AV-2014-069, 2014.

\bibitem{eu261}
European Parliament and Council, Regulation~(EC) No~261/2004 establishing common rules on compensation and assistance to passengers in the event of denied boarding and of cancellation or long delay of flights, February~2004.

\bibitem{peterson2021diversion}
D.~C.~Peterson et~al., ``Data-Driven Estimation of the Impact of Diversions Due to In-Flight Medical Emergencies on Flight Delay and Aircraft Operating Costs,'' \emph{Aerospace Medicine and Human Performance}, vol.~92, no.~2, pp.~81--87, 2021.
\end{thebibliography}

\subsection{描述性结果}
图 \ref{fig:monthly-trend-zh} 展示了 2020 至 2025 年 JFK 月度延误和取消数量的变化。最高延误月份出现在 2023 年 7 月，当月 15 分钟以上延误达到 3,657 次，延误率为 32.69\%。最高取消月份则出现在 2020 年 3 月，当月取消 1,900 次，反映出疫情初期对航空运行的特殊冲击。这个结果提醒我们，延误与取消并不总是同步变化。延误更多体现运行压力和拥堵，取消则更容易受到外部突发事件和大规模计划调整影响。

\begin{figure}[H]
\centering
\includegraphics[width=0.94\linewidth]{chart_1_monthly_trend.png}
\caption{JFK 月度延误与取消趋势}
\label{fig:monthly-trend-zh}
\end{figure}

从原因结构看，Airline 和 Late Aircraft 是最主要的延误分钟来源，分别占总延误分钟的 36.16\% 和 35.14\%。NAS 占 24.32\%，Weather 与 Security 合计不足 5\%。这与前文的风险识别形成呼应：天气虽然是航空风险叙事中非常直观的外部因素，但在该样本中，日常风险的主体仍然来自内部运行链条和系统拥堵。

\begin{figure}[H]
\centering
\includegraphics[width=0.92\linewidth]{chart_2_delay_cause_breakdown.png}
\caption{延误原因的分钟贡献结构}
\label{fig:cause-breakdown-zh}
\end{figure}

\begin{table}[H]
\centering
\caption{主要延误原因占比}
\begin{tabular}{lrr}
\toprule
原因 & 延误分钟 & 占比 \\
\midrule
Airline & 4,022,961 & 36.16\% \\
Late Aircraft & 3,910,112 & 35.14\% \\
NAS & 2,705,624 & 24.32\% \\
Weather & 466,088 & 4.19\% \\
Security & 21,053 & 0.19\% \\
\bottomrule
\end{tabular}
\end{table}

图 \ref{fig:airline-profile-zh} 将航司层面的延误率、取消率与平均延误严重度放在同一张风险画像中。JetBlue Airways 的总延误分钟最高，Hawaiian Airlines 的延误率最高，而 Envoy Air 的平均延误严重度最高。这样的结果不应被解释为简单的“哪家航司最差”，因为不同航司的航班量、航线结构和样本规模不同。它更适合作为后续模型的分层视角，帮助我们识别哪些航司月份可能对尾部风险贡献更大。

\begin{figure}[H]
\centering
\includegraphics[width=0.88\linewidth]{chart_3_airline_risk_profile.png}
\caption{航司风险画像}
\label{fig:airline-profile-zh}
\end{figure}

\subsection{对后续模型构建的帮助}
这一部分的数据工作对后续模型有三方面帮助。第一，清洗后的核心数据把频率变量和严重度变量分开，使后续可以分别拟合 Poisson、Negative Binomial、Lognormal 或 Weibull 等候选分布。第二，延误原因汇总把 BTS 的数据库口径转化为课程中的运营风险语言，后续可以把 Airline 与 Late Aircraft 合并为内部运营链条风险，把 NAS 视为系统约束风险，把 Weather 与 Security 视为外部冲击风险。第三，月度汇总保留了连续时间顺序，后续可以进一步定义 normal 与 disrupted 两类运行状态，并估计状态转移关系。

因此，本章节的核心贡献不是某一个单独图表，而是把原始航班运行数据整理成一套可解释、可复现、可建模的风险数据底座。它让前文的管理分析能够落到量化指标上，也让后文的统计模型不只是形式上的分布拟合，而是有清楚业务含义的运营风险分析。

\subsection{局限与下一步}
本阶段仍有两点局限。第一，数据粒度是航司月份层面，不能观察单个航班、机场时段或具体天气事件，因此短时拥堵与航班链式传播只能通过月度汇总间接体现。第二，延误原因由 BTS 按规则分摊，部分原因计数字段可能出现小数，这意味着它们更像风险归因指标，而不是严格的独立事件计数。后续建模在使用这些变量时，需要明确说明这一口径。

尽管存在这些限制，当前数据集已经足以支撑课程项目要求中的定量分析部分。后续章节可以在本章输出的基础上，继续完成频率分布比较、严重度分布比较、状态依赖建模和 aggregate risk 情景模拟。

\end{document}
